\chapter{The Which-Strategy-When Decision Framework}
\label{ch:which_strategy_when}
\chaptermark{Which strategy when}

Part V gave us an arsenal of strategies
(\cref{ch:market_making,ch:mean_reversion_ou,ch:pairs_cointegration,ch:basket_etf_arb,ch:momentum,ch:arbitrage_zoo,ch:greeks_hedging,ch:vol_surface,ch:almgren_chriss,ch:kelly_risk,ch:ml_trading}),
each taught inside its own assumed market conditions. This chapter
answers the practitioner question: given the state of the market
\emph{right now}, which of those strategies should run, on which
product, with what parameters, and which proposed trades should
fire? The output is a concrete list of sized, scheduled trades.

The machinery has four layers, assembled from \citeChanQT, the CPO
blog \citep{chan2020cpo}, the QuantStart HMM tutorial
\citep{quantstart2016hmm}, and \citeAFML. A \textbf{Hidden Markov
Model} (HMM) detects regimes on the slow clock (regimes last
weeks). A \textbf{decision tree} maps observables to a strategy on
the strategy's own clock (per bar/tick). Chan's \textbf{Conditional
Parameter Optimisation} (CPO) adapts parameters; its regression
is retrained nightly/weekly but queried per tick.
\textbf{Meta-labeling} (\cref{ch:ml_trading}) runs fastest, once
per proposed trade. Downstream, Kelly sizes and Almgren--Chriss
schedules. Thinking about each layer's clock is half the
engineering; making each layer's output legible to the next is the
other half.

\begin{backgroundbox}[Prerequisites]
From earlier in the book:
\begin{itemize}
  \item \Cref{ch:mean_reversion_ou}: OU process, half-life
  $\tau_{1/2} = \ln 2 / \theta$, ADF test, $z$-score entry rule
  (the mean-reversion branch is gated on ADF $p$ and $z$).
  \item \Cref{ch:pairs_cointegration,ch:basket_etf_arb}:
  cointegration, basket-vs-constituents spreads, Engle--Granger
  (basket branch gated on basis $z$).
  \item \Cref{ch:momentum}: TSMOM, cross-sectional momentum,
  12-month look-back (trend branch gated on sign of trailing
  12-month return).
  \item \Cref{ch:market_making}: Avellaneda--Stoikov; market
  making is net-profitable when spread exceeds per-fill
  fair-value uncertainty and inventory drift is manageable.
  \item \Cref{ch:greeks_hedging,ch:vol_surface}: long gamma
  makes money when realised vol exceeds implied.
  \item \Cref{ch:ml_trading}: triple-barrier labels and
  meta-labeling, the trade-level selector in the framework
  below.
  \item \Cref{ch:backtesting}: walk-forward and CPCV. Every
  component must be validated under these rules.
  \item \Cref{ch:kelly_risk}: Kelly sizing.
\end{itemize}
From outside the book: elementary HMM vocabulary (hidden states,
transition matrix, forward/backward, EM; Rabiner's 1989 tutorial
is the standard reference); basic regression and random forests at
the level of \cref{ch:ml_trading}.
\end{backgroundbox}

% =====================================================================
\section{A concrete morning-of-trading decision}
\label{sec:wsw_toy_example}

Before the formal machinery, walk through a toy morning at a
Prosperity-style desk with four products on the screen. This is how
the framework should \emph{feel}; the rest of the chapter explains
how each step is computed. The universe:
\begin{itemize}
  \item \texttt{RAINFOREST\_RESIN}: nearly flat price around
  $\approx 10{,}000$ seashells, with a tight, stable spread. No
  information-driven flow. A textbook candidate for passive
  market making (\cref{ch:market_making}).
  \item \texttt{KELP}: mid-volatility, mean-reverting on the
  minute-scale time axis. Historical ADF $p$-value small, OU
  half-life on the order of tens of ticks. A textbook
  $z$-score mean-reversion candidate
  (\cref{ch:mean_reversion_ou}).
  \item \texttt{SQUID\_INK}: volatile, occasionally driven by
  an informed trader (Olivia in the Prosperity 3 worked
  examples of \cref{ch:momentum}). When she is buying, the
  price runs in the same direction for many ticks. A textbook
  TSMOM candidate, conditional on the informed-flow signal.
  \item \texttt{PICNIC\_BASKET1}: a known basket of
  individual products (6 \texttt{CROISSANTS}, 3
  \texttt{JAMS}, 1 \texttt{DJEMBE}). The basket and its
  replicating portfolio can deviate from each other; the
  spread is a classic basket-arb candidate
  (\cref{ch:basket_etf_arb}).
\end{itemize}

At 08:00 local, a single tick of diagnostic statistics is computed
from the last five minutes of order-book data. We show them in
\cref{tab:wsw_toy_observables}. The rows are observables; the
columns are products.

\begin{table}[htbp]
\centering
\begin{tabular}{lrrrr}
\toprule
Observable ($o_t$ component)
  & \texttt{RESIN} & \texttt{KELP} & \texttt{SQUID}
  & \texttt{BASKET1} \\
\midrule
5-min realised vol (bps) & 1.8 & 12.0 & 38.0 & 14.0 \\
Quoted half-spread (bps) & 0.5 & 1.5 & 4.0 & 6.0 \\
Top-of-book depth (units) & 40 & 30 & 15 & 20 \\
Lag-1 price autocorr & $+0.01$ & $-0.35$ & $+0.18$ & $-0.02$ \\
ADF $p$-value (last 500 ticks) & 0.02 & 0.005 & 0.41 & n/a \\
12-min trailing return ($\Delta p / p$) & $+0.0001$ & $+0.0004$ & $+0.0120$ & $-0.0003$ \\
Basket basis $z$-score & n/a & n/a & n/a & $+2.1$ \\
Olivia (informed) buy flag & no & no & \textbf{yes} & no \\
\bottomrule
\end{tabular}
\caption{Morning-of-trading diagnostic statistics for the toy
universe of \cref{sec:wsw_toy_example}. All four strategies in
the pool are discriminable from these eight numbers.}
\label{tab:wsw_toy_observables}
\end{table}

Reading \cref{tab:wsw_toy_observables} row by row: \texttt{RESIN}
has tiny vol (1.8~bps), a non-trivial 0.5~bps half-spread (so full
spread $\approx$ half the per-5-min vol), and zero lag-1
autocorrelation: no trend, no mean-reverting edge. Passive
quoting is the only positive-expectancy move. \textbf{Select:
market making}.

\texttt{KELP} has order-of-magnitude larger vol, lag-1
autocorrelation $-0.35$ (mean-reverting), ADF $p = 0.005$
(rejecting unit root), and half-spread small relative to vol. Assume
$|z| > 1.5$ (otherwise no trade fires). \textbf{Select: OU
$z$-score mean reversion}.

\texttt{SQUID\_INK} has the Olivia flag on, $+1.2\%$ trailing
return, positive lag-1 autocorrelation, and ADF $p = 0.41$ (unit
root, so no mean-reversion branch fires). Trend is present and
informed-trader-confirmed. \textbf{Select: time-series momentum
long}.

For \texttt{PICNIC\_BASKET1} the tradeable is the basket spread
(\cref{ch:basket_etf_arb}). Basis $z = +2.1$ exceeds the $+1.5$
entry threshold: basket is rich relative to replicating portfolio.
\textbf{Select: short basket / long replicating portfolio}.

Four products, eight numbers, four strategies: each selection is
the gating condition of a previously-written chapter, read off a
single table. The rest of this chapter formalises this.

\begin{prosperitybox}
\textbf{Prosperity 3: each selected strategy is on the leaderboard.}
All four strategies in \cref{sec:wsw_toy_example} appear in the
Frankfurt Hedgehogs' second-place write-up: market making on flat
low-vol products, OU mean reversion on mid-vol mean-reverting
products, basket arb on \texttt{PICNIC\_BASKET1}/\texttt{2}, and
TSMOM on \texttt{SQUID\_INK} conditional on the Olivia signal
(seven strategies combined). This chapter formalises ``how did they
pick which strategy to run where?'' from \citeChanQT{} and
\citep{chan2020cpo}; the write-up is what the output looks like
when done right.
\end{prosperitybox}

% =====================================================================
\section{Regime detection via Hidden Markov Models}
\label{sec:wsw_hmm}

The toy example cheated: observables drift with the \emph{market
regime}, the hidden state that makes returns, spreads,
volatility, and correlations covary. A $z = 2$ in a calm regime is
real; in a crisis regime $z = 2$ is small relative to intraday
shocks and means something different. A \textbf{regime-aware}
framework wraps the decision tree in a layer that first estimates
the regime, then reads the decision tree under that regime's
conditional distribution.

We model the market as a \textbf{Hidden Markov Model} (HMM),
following \citep{rabiner1989tutorial} (the notation used by
\texttt{hmmlearn}).

\begin{definition}[Hidden Markov Model]
\index{hidden Markov model (HMM)|textbf}\index{HMM|textbf}
\label{def:wsw_hmm}
A Hidden Markov Model with $N$ states is a triple
$(\boldsymbol{\pi}, A, B)$ where
\begin{itemize}
  \item $\boldsymbol{\pi} = (\pi_1, \ldots, \pi_N)$ is the
  initial state distribution, $\pi_i = \Prob(S_1 = i)$.
  \item $A \in \R^{N\times N}$ is the transition matrix,
  $A_{ij} = \Prob(S_{t+1} = j \mid S_t = i)$.
  \item $B = \{b_i(\cdot)\}_{i=1}^N$ is the set of
  emission densities: conditional on being in state $i$ at
  time $t$, the observation $o_t$ is drawn from $b_i$.
\end{itemize}
In the trading application, the observation $o_t$ is typically a
daily or intraday return (or a vector of features) and each
emission $b_i$ is a Gaussian (or a Gaussian mixture) over that
return. The hidden state sequence $S_1, S_2, \ldots, S_T$ is
unobserved; all one sees is $o_1, \ldots, o_T$.
\end{definition}

Given the model, three algorithms do all the work. We state them in
one place and refer back.

\begin{keyfactbox}[The three HMM algorithms]
Given an HMM $(\boldsymbol{\pi}, A, B)$ and an observation sequence
$o_{1:T}$:
\begin{enumerate}
  \item \textbf{Forward algorithm.} Computes the filtered
  posterior $\Prob(S_t = i \mid o_{1:t})$ for each state $i$
  and time $t$, in $O(N^2 T)$ operations. Gives you ``what
  regime is the market \emph{probably} in right now, given
  everything I have seen so far.''
  \item \textbf{Viterbi algorithm.} Computes the most likely
  single state path
  $\arg\max_{S_{1:T}} \Prob(S_{1:T} \mid o_{1:T})$, in
  $O(N^2 T)$ operations. Gives you a single hard assignment
  of a regime label to every historical time, useful for
  backtesting regime-conditional returns.
  \item \textbf{Baum--Welch (EM).} Estimates the parameters
  $(\boldsymbol{\pi}, A, B)$ from data by expectation
  maximisation, using forward--backward as the E-step. Gives
  you a \emph{trained} HMM in the first place.
\end{enumerate}
For the regime-detection use case we fit parameters once with
Baum--Welch on a long training window, then at every decision
tick we run the forward algorithm on the recent observation
history and read the current-state posterior off the last column
\citep{quantstart2016hmm}.
\end{keyfactbox}

We do not re-derive the three algorithms; the derivation is
four pages of index algebra that reduces to ``law of total
probability plus memoisation'' (see \citep{rabiner1989tutorial}).
Define the forward variable
\begin{equation}
  \alpha_t(i) \;\defeq\; \Prob(o_1, o_2, \ldots, o_t,\, S_t = i),
  \label{eq:wsw_forward_var}
\end{equation}
the joint probability of having seen the observations so far and
the hidden state at the current time being $i$. Its recursion is
\begin{equation}
  \alpha_{t}(j)
  \;=\;
  \Big[\sum_{i=1}^{N} \alpha_{t-1}(i)\, A_{ij}\Big]\, b_j(o_t),
  \qquad \alpha_1(j) = \pi_j\, b_j(o_1),
  \label{eq:wsw_forward_recursion}
\end{equation}
which is the law of total probability: arrive in any $i$ at $t-1$,
transition to $j$, emit $o_t$. Summing $\alpha_T$ over states
gives the total likelihood $\Prob(o_{1:T})$ as a by-product. The
filtered posterior is
\begin{equation}
  \Prob(S_t = i \mid o_{1:t})
  \;=\;
  \frac{\alpha_t(i)}{\sum_{j=1}^{N}\alpha_t(j)}.
  \label{eq:wsw_forward_posterior}
\end{equation}
In practice $\alpha_t$ underflows rapidly (each step multiplies by
an emission density $\sim 30$), so the recursion is re-normalised
every step: divide $\alpha_t(\cdot)$ by its sum, making it the
posterior directly. This is the form used below and in every
production implementation.

\subsection{A 2-state HMM on synthetic SPY-style returns}
\label{subsec:wsw_hmm_worked}

Our model has $N = 2$ hidden states. State $S = 1$ is ``low-vol
bull,'' state $S = 2$ is ``high-vol bear.'' Emission distributions
are Gaussian with parameters
\begin{equation}
  (\mu_1, \sigma_1) = (+0.05\%,\; 0.80\%),
  \qquad
  (\mu_2, \sigma_2) = (-0.10\%,\; 2.00\%),
  \label{eq:wsw_hmm_emissions}
\end{equation}
per day (daily returns, not log-prices), representative of SPY
over 1993--2014 \citep{quantstart2016hmm}. Sticky transition
matrix,
\begin{equation}
  A \;=\;
  \begin{pmatrix}
    0.98 & 0.02 \\
    0.02 & 0.98
  \end{pmatrix},
  \label{eq:wsw_hmm_A}
\end{equation}
giving mean regime dwell time $1/0.02 = 50$ days (two to three
trading months). Initial distribution $\boldsymbol{\pi} = (0.5,
0.5)$. Observation sequence:
\begin{equation}
  o_{1:5} \;=\; (+0.40\%,\; +0.60\%,\; -0.30\%,\; -2.50\%,\; -3.00\%).
  \label{eq:wsw_hmm_obs}
\end{equation}
Mild-up followed by a sharp two-day drawdown. We run the forward
algorithm step by step, reporting the normalised posterior. All
digits are reproducible against \texttt{scipy.stats.norm.pdf}.

\paragraph{Step 1: $t = 1$, $o_1 = +0.40\%$.} Emission likelihoods:
\begin{align}
  b_1(o_1)
  &= \Normal(0.0040 \mid 0.0005,\, 0.0080)
  = 45.32, \\
  b_2(o_1)
  &= \Normal(0.0040 \mid -0.0010,\, 0.0200)
  = 19.33.
\end{align}
State 1's density is just over twice state 2's (observation is
inside state 1's 1$\sigma$ band, at the edge of state 2's).
Multiply by prior and renormalise:
\begin{equation}
  \alpha_1 \;=\;
  \frac{(0.5)(45.32)}{(0.5)(45.32) + (0.5)(19.33)}
  \;=\;
  \underbrace{0.7010}_{\Prob(S_1 = 1 \mid o_1)}.
  \label{eq:wsw_hmm_alpha1}
\end{equation}

\paragraph{Step 2: $t = 2$, $o_2 = +0.60\%$.} The recursion is
\begin{equation}
  \alpha_t(j)
  \;=\;
  \Big[\sum_{i=1}^{N} \alpha_{t-1}(i)\, A_{ij}\Big]\,
  \underbrace{b_j(o_t)}_{\text{emission}},
  \label{eq:wsw_hmm_recursion}
\end{equation}
with re-normalisation after each step. Emission likelihoods at
$o_2 = 0.006$ are $b_1 = 39.37$, $b_2 = 18.76$. Propagate:
\begin{align}
  \alpha_2(1)
  &\propto (0.7010\cdot 0.98 + 0.2990\cdot 0.02)\,b_1(0.006)
  \;=\; 0.6929 \cdot 39.37
  \;=\; 27.28,\\
  \alpha_2(2)
  &\propto (0.7010\cdot 0.02 + 0.2990\cdot 0.98)\,b_2(0.006)
  \;=\; 0.3071 \cdot 18.76
  \;=\; 5.76.
\end{align}
After renormalisation, $\alpha_2 = (0.8256,\, 0.1744)$: a second
low-vol-consistent observation sharpens the inference.

\paragraph{Step 3: $t = 3$, $o_3 = -0.30\%$.} Still inside state
1's 1$\sigma$ band: $\alpha_3 = (0.9083,\, 0.0917)$.

\paragraph{Step 4: $t = 4$, $o_4 = -2.50\%$.} The model changes its
mind. $-2.5\%$ is $-3.1\sigma$ under state 1 but only $-1.25\sigma$
under state 2. Emission ratio:
\begin{equation}
  \frac{b_2(o_4)}{b_1(o_4)}
  \;=\;
  \frac{\Normal(-0.025 \mid -0.001,\, 0.020)}
       {\Normal(-0.025 \mid +0.0005,\, 0.008)}
  \;=\; \frac{9.71}{0.31}
  \;=\; 31.3.
  \label{eq:wsw_hmm_flip}
\end{equation}
State 2's emission is 31$\times$ more likely. The sticky $0.98$
diagonal cannot hold state 1 against this:
\begin{align}
  \alpha_4(1)
  &\propto (0.9083\cdot 0.98 + 0.0917\cdot 0.02)(0.31)
  \;=\; 0.2772, \\
  \alpha_4(2)
  &\propto (0.9083\cdot 0.02 + 0.0917\cdot 0.98)(9.71)
  \;=\; 1.0514.
\end{align}
After renormalisation, $\alpha_4 = (0.2086,\, 0.7914)$. The posterior
flipped onto state 2 in one observation.

\paragraph{Step 5: $t = 5$, $o_5 = -3.00\%$.} A second high-vol
observation locks it in: $\alpha_5 = (0.0014,\, 0.9986)$.

\begin{table}[htbp]
\centering
\begin{tabular}{crrr}
\toprule
$t$ & $o_t$ & $\Prob(S_t = 1 \mid o_{1:t})$
  & $\Prob(S_t = 2 \mid o_{1:t})$ \\
\midrule
1 & $+0.40\%$ & $\mathbf{0.7010}$ & $0.2990$ \\
2 & $+0.60\%$ & $\mathbf{0.8256}$ & $0.1744$ \\
3 & $-0.30\%$ & $\mathbf{0.9083}$ & $0.0917$ \\
4 & $-2.50\%$ & $0.2086$ & $\mathbf{0.7914}$ \\
5 & $-3.00\%$ & $0.0014$ & $\mathbf{0.9986}$ \\
\bottomrule
\end{tabular}
\caption{Forward-algorithm posteriors for the 2-state HMM of
\cref{subsec:wsw_hmm_worked}. The regime flip happens in a single
observation at $t = 4$, when the emission likelihood of the
high-vol state exceeds that of the low-vol state by a factor of
31. Bold entries mark the argmax state at each time.}
\label{tab:wsw_hmm_alphas}
\end{table}

Two takeaways. The sticky matrix prevents overreaction to small
observations ($-0.3\%$ did not flip it) but large observations
overwhelm it in one step ($-2.5\%$ did). And once flipped, the
model is confident: day 5 pushes to $0.9986$ because the ratio acts
on an already-tilted posterior.

\begin{intuitionbox}
An HMM detects regimes by racing two densities. The posterior is
the previous posterior (propagated through the transition matrix)
times the new observation's emission likelihoods. Flat observations
preserve it; observations deep in one state's tail and inside
another's body snap it over. Stickiness controls the evidence
required: large off-diagonals flip on one bad day, small ones
require several. Choosing stickiness is the single most important
HMM parameter decision. Too nervous, and the pipeline swaps
strategies daily; too confident, and it sits in the wrong regime for
a month.
\end{intuitionbox}

\begin{pitfallbox}[Number of states and the Gaussianity assumption]
Two failure modes dominate. \textbf{Too many states}: $N = 2$ finds
``quiet vs crisis''; $N = 3$ sometimes finds ``up/range/down'',
sometimes garbage; $N = 6$ almost always overfits: states
capture one or two outliers and the posterior flickers. Start at
$N = 2$; only increase with a-priori reason and out-of-sample
evidence. \textbf{Gaussian emissions on raw returns}: daily equity
returns are fat-tailed, so the Gaussian HMM assigns tails to a
phantom high-vol state. Use Gaussian mixtures per state or fit on
GARCH-standardised residuals. \citeChanQT{} and
\citep{quantstart2016hmm} both recommend small $N$ and careful
residualisation.
\end{pitfallbox}

% =====================================================================
\section{What regimes mean for strategy selection}
\label{sec:wsw_regime_map}

Given a regime, which strategies should run? The stylised facts
from \citeChan{}, \citeChanQT{}, and the regime-conditional
backtest literature:
\begin{itemize}
  \item \textbf{Low-vol bull}: tight spreads, thick books, slow
  upward drift. TSMOM (\cref{ch:momentum}) earns its highest
  Sharpe. Mean reversion is anti-correlated with the trend and
  gets run over. Market making earns a steady spread (low adverse
  selection). Long gamma loses: realised vol stays below implied.
  \item \textbf{High-vol bear}: wide spreads, thin books. TSMOM
  short side earns highest Sharpe (2008 is the textbook case);
  cross-sectional momentum still works at reduced Sharpe. Mean
  reversion catches falling knives. Market making is dangerous
  (the book walks); \citeChan{} recommends aggressive inventory
  skew. Long gamma earns its highest returns.
  \item \textbf{Range-bound}: tightly bracketed, tight spreads,
  short half-lives. OU mean reversion
  (\cref{ch:mean_reversion_ou}) and pairs
  (\cref{ch:pairs_cointegration}) earn their highest Sharpes.
  TSMOM whipsaws. Market making earns the spread. Basket arb is
  a function of basis dislocation only, regime-independent.
  \item \textbf{Regime transition}: the worst place. HMM
  posterior uncertain, neither trend nor mean-reversion reliable,
  execution costs elevated. Every practitioner guide recommends
  \emph{reducing position size} rather than trading through.
\end{itemize}

\Cref{tab:wsw_regime_map} summarises. Entries are qualitative
relative-Sharpe assessments from \citeChan{,}\citeChanQT{}; the
\emph{sign} of each cell is robust across decades of backtests, the
magnitude is not.

\begin{table}[htbp]
\centering
\begin{tabular}{lcccc}
\toprule
Strategy family
  & Low-vol bull & High-vol bear
  & Range-bound & Transition \\
\midrule
Market making
  & $+$ & $-$ & $++$ & $-$ \\
OU / pairs mean reversion
  & $-$ & $--$ & $++$ & $-$ \\
Basket arbitrage
  & $\circ$ & $\circ$ & $\circ$ & $\circ$ \\
TSMOM (long only)
  & $++$ & $-$ & $-$ & $-$ \\
TSMOM (long/short)
  & $++$ & $+$ & $-$ & $-$ \\
Long gamma / long vol
  & $-$ & $++$ & $-$ & $+$ \\
Arbitrage sweep
  & $+$ & $+$ & $+$ & $+$ \\
\bottomrule
\end{tabular}
\caption{Qualitative regime-to-strategy map. Cells are
$++$ (strongest Sharpe), $+$ (positive), $\circ$ (regime-neutral,
function of basis only), $-$ (loses money), $--$ (loses money
fast). Basket arbitrage and the arbitrage sweep are regime-neutral
because they exploit a price relationship, not a direction.
Synthesised from \citeChan{}, \citeChanQT{}, and \citep{quantstart2016hmm}.}
\label{tab:wsw_regime_map}
\end{table}

The specific cells matter less than the pattern: most strategies have
\emph{sharp sign flips} across regimes. Running mean reversion
through a bull trend, or TSMOM through a range, is the most common
money-loss pattern.

Three rows deserve a closer look. \textbf{Market making} is $+$ in
low-vol bull but $++$ in range-bound. The reason is adverse
selection (\cref{ch:kyle_gm}): in a trend, informed flow picks off
the passive quote on one side; in a range, both sides come from
uninformed noise traders so the spread is earned on every
round-trip. Market making profits when adverse selection is small
relative to quoted spread, and that is smallest in range-bound
regimes.

\textbf{Long gamma} is $+$ in transition despite being $-$ in
low-vol. Long gamma makes money when realised vol exceeds the
implied at purchase, and regime transitions are precisely when
implied hasn't caught up to the new regime's realised. ``Vol trades
buy dislocations.''

\textbf{Basket arbitrage} is $\circ$ (regime-neutral) in every
column because the basket spread is a price \emph{relationship},
not a direction. A trend in both legs leaves the spread unchanged.
The basket branch comes \emph{before} regime-conditional branches
for this reason. On the Frankfurt Hedgehogs' leaderboard basket arb
P\&L was top-three: ``regime-neutral'' means ``runs in every
regime,'' not ``small.''

\begin{intuitionbox}
One-sentence summary: \emph{strategies that bet on price
\textbf{direction} have sharp sign flips across regimes;
strategies that bet on a \textbf{relationship} between prices do
not}. Momentum and mean reversion are directional bets on
autocorrelation. Market making is a directional bet on the sign of
adverse selection. Arbitrage and basket spreads are relationship
bets: they pay when the relationship is dislocated regardless
of leg direction. Long gamma is realised-vs-implied (a
relationship), less regime-sensitive than directional strategies
but more than cross-sectional arbitrage because the relationship
itself moves with regime.
\end{intuitionbox}

A synthetic check makes the sign flip concrete. On 2000-step
series, half trending (drift $\pm 0.001$, $\sigma = 0.005$) and
half mean-reverting (OU with $\theta = 0.1$, $\sigma = 0.3$), a
20-step TSMOM and a 50-step $|z|>1.5$ OU strategy gave:
\begin{itemize}
  \item Trending regime: TSMOM Sharpe $= +1.65$,
  mean-reversion Sharpe $= -1.64$.
  \item Mean-reverting regime: TSMOM Sharpe $= -1.97$,
  mean-reversion Sharpe $= +2.22$.
\end{itemize}
Each strategy is $\approx +2$ in its regime and $\approx -2$ in
the other. A regime-agnostic trader running both nets $\approx 0$
and eats all costs as drag; a regime-aware one realises $+2$ in
both. That is the framework, in four numbers.

% =====================================================================
\section{The opportunity-diagnosis decision tree}
\label{sec:wsw_decision_tree}

\begin{figure}[htbp]
\centering
\begin{tikzpicture}[
  node distance=6mm and 10mm,
  every node/.style={font=\footnotesize},
  guard/.style={draw, rectangle, rounded corners=2pt,
    fill=blue!5, align=center, minimum width=40mm,
    minimum height=7mm, inner sep=3pt},
  fire/.style={draw, rectangle, fill=green!10, align=center,
    minimum width=22mm, minimum height=5mm, inner sep=2pt},
  skip/.style={draw, rectangle, fill=gray!10, align=center,
    minimum width=22mm, minimum height=5mm, inner sep=2pt},
  arrow/.style={-Stealth, thick}
]
  \node[guard] (g1) {Guard 1: classical no-risk arb available?};
  \node[fire, right=of g1] (f1) {\texttt{ARB}};
  \node[guard, below=of g1] (g2) {Guard 2: basket, $|z^{(\text{basis})}|>1.5$?};
  \node[fire, right=of g2] (f2) {\texttt{BASKET}};
  \node[guard, below=of g2] (g3) {Guard 3: ADF-$p<0.05$, $|z|>z_{\text{entry}}$?};
  \node[fire, right=of g3] (f3) {\texttt{OU}/\texttt{PAIRS}};
  \node[guard, below=of g3] (g4) {Guard 4: $|r^{(12\text{m})}|>\theta_{\text{trend}}$ \& trending?};
  \node[fire, right=of g4] (f4) {\texttt{TSMOM}};
  \node[guard, below=of g4] (g5) {Guard 5: $\hat\sigma>\sigma^{\text{impl}}$, gamma tradable?};
  \node[fire, right=of g5] (f5) {\texttt{GAMMA}};
  \node[guard, below=of g5] (g6) {Guard 6: $\hat s > 2\hat\sigma/\sqrt{N_f}$, depth OK?};
  \node[fire, right=of g6] (f6) {\texttt{MM}};
  \node[skip, below=of g6] (gskip) {\texttt{SKIP}};

  \draw[arrow] (g1) -- node[above,font=\tiny]{yes} (f1);
  \draw[arrow] (g2) -- node[above,font=\tiny]{yes} (f2);
  \draw[arrow] (g3) -- node[above,font=\tiny]{yes} (f3);
  \draw[arrow] (g4) -- node[above,font=\tiny]{yes} (f4);
  \draw[arrow] (g5) -- node[above,font=\tiny]{yes} (f5);
  \draw[arrow] (g6) -- node[above,font=\tiny]{yes} (f6);
  \draw[arrow] (g1) -- node[left,font=\tiny]{no} (g2);
  \draw[arrow] (g2) -- node[left,font=\tiny]{no} (g3);
  \draw[arrow] (g3) -- node[left,font=\tiny]{no} (g4);
  \draw[arrow] (g4) -- node[left,font=\tiny]{no} (g5);
  \draw[arrow] (g5) -- node[left,font=\tiny]{no} (g6);
  \draw[arrow] (g6) -- node[left,font=\tiny]{no} (gskip);
\end{tikzpicture}
\caption{The opportunity-diagnosis decision tree of
\cref{sec:wsw_decision_tree}, drawn as a priority-ordered guard
chain. Each guard is a threshold on one component of the
observable vector $o_t^{(p)}$. The first guard whose condition
fires wins; if no guard fires, the tree outputs \texttt{SKIP}.
Each thresholded parameter ($z_{\text{entry}}$,
$\theta_{\text{trend}}$, the depth cut, $\sigma^{\text{impl}}$)
is regime-conditional through CPO (\cref{sec:wsw_cpo}).}
\label{fig:wsw_decision_tree}
\end{figure}


Generalise the toy table. Define the \textbf{observable vector} at
time $t$ for product $p$ as
\begin{equation}
  o_t^{(p)} \;=\;
  \big(
    \hat{\sigma}_t,\,
    \hat{s}_t,\,
    \hat{d}_t,\,
    \rho_t^{(1)},\,\rho_t^{(5)},\,
    \widehat{\mathrm{OFI}}_t,\,
    z_t^{(\mathrm{basis})},\,
    \text{ADF-}p_t,\,
    r_t^{(12\mathrm{m})},\,
    \mathrm{flags}_t
  \big),
  \label{eq:wsw_obs_vec}
\end{equation}
where $\hat{\sigma}_t$ is realised vol, $\hat{s}_t$ half-spread,
$\hat{d}_t$ top-of-book depth, $\rho_t^{(k)}$ lag-$k$ autocorr,
$\widehat{\mathrm{OFI}}_t$ order-flow imbalance
(\cref{ch:order_flow}), $z_t^{(\mathrm{basis})}$ basket basis
$z$-score, $\text{ADF-}p_t$ stationarity $p$-value,
$r_t^{(12\mathrm{m})}$ trailing 12-month return, and
$\mathrm{flags}_t$ categorical state (informed-trader, earnings,
expiry).

The \textbf{strategy pool} $\mathcal{S}$: market making
(\texttt{MM}), OU $z$-score (\texttt{OU}), cointegration pairs
(\texttt{PAIRS}), basket arbitrage (\texttt{BASKET}), TSMOM
(\texttt{TSMOM}), long-vol gamma (\texttt{GAMMA}), catch-all
arbitrage (\texttt{ARB}) for put-call parity, triangular, and
locational.

The decision tree $\mathcal{T}: o_t^{(p)} \to \mathcal{S} \cup
\{\text{SKIP}\}$ is a priority chain of guards: each guard is a
threshold on one component of $o_t^{(p)}$, and the first that
fires wins. Arbitrage comes first (near-free money), then basket
arbitrage (regime-neutral), then regime-conditional strategies.

\begin{keyfactbox}[Decision-tree scaffold for one product]
At each decision tick, for each product $p$ in the universe:
\begin{enumerate}
  \item If any classical no-risk arbitrage (put--call parity,
  triangular, locational) is available with $P\&L > c_{\text{exec}}$,
  fire \texttt{ARB}. Continue to next product.
  \item Else, if $p$ is a basket and
  $|z_t^{(\mathrm{basis})}| > 1.5$, fire \texttt{BASKET} in the
  direction opposite the basis dislocation.
  \item Else, if ADF-$p < 0.05$ and
  $|z_t^{(p)}| > z_{\text{entry}}(\text{regime})$, fire
  \texttt{OU} (or \texttt{PAIRS} if $p$ is a pair). Both bet on
  reversion to a stationary mean; \texttt{OU} uses the product's
  own price, \texttt{PAIRS} uses the cointegrated spread between
  two products (\cref{ch:pairs_cointegration}).
  \item Else, if $|r_t^{(12\mathrm{m})}| > \theta_{\text{trend}}$
  and the HMM regime is trending, fire \texttt{TSMOM}.
  \item Else, if $\hat{\sigma}_t > \sigma^{\mathrm{impl}}$
  (realised exceeds implied) and gamma is tradable, fire
  \texttt{GAMMA}.
  \item Else, if $\hat{s}_t > 2 \hat{\sigma}_t/\sqrt{N_{\text{fills}}}$
  and $\hat{d}_t$ is above a minimum depth, fire \texttt{MM}.
  The inequality says the quoted half-spread exceeds the
  per-fill fair-value uncertainty, so each round-trip is
  expected-positive after adverse selection.
  \item Else, \texttt{SKIP} (no strategy fires for this product
  at this tick).
\end{enumerate}
Each threshold is a tunable parameter; in the next section,
CPO makes them all regime-conditional.
\end{keyfactbox}

Guards are ordered by \emph{priority}, not likelihood. Arbitrage
first (regime-free, riskless). Basket arbitrage next
(regime-neutral, \cref{tab:wsw_regime_map}). Mean reversion and
TSMOM in the middle, mutually exclusive via the
cointegration/regime tests. Market making near the bottom (default
when no stronger signal). SKIP as last resort, more often
correct than beginners think. The most common money-loss in
Prosperity 3 write-ups was strategies firing every tick because the
author refused to hold no position.

\subsection{Worked guard evaluation}
\label{subsec:wsw_guard_eval}

Replay the toy table guard-by-guard with arithmetic.

\paragraph{Guard 1 (arbitrage sweep).} No options, fiat pairs, or
multi-venue listings in the toy universe; does not fire.

\paragraph{Guard 2 (basket arbitrage).} Only \texttt{PICNIC\_BASKET1}
is a basket; basis $z = +2.1 > 1.5$. Fires: short \texttt{BASKET1},
long $6 \cdot$\texttt{CROISSANTS} $+ 3 \cdot$\texttt{JAMS} $+ 1
\cdot$\texttt{DJEMBE}. Route and continue.

\paragraph{Guard 3 (OU mean reversion).} Check ADF-$p < 0.05$ and
$|z| > z_{\text{entry}}$. \texttt{RESIN}: ADF-$p = 0.02$ but
autocorr $+0.01$ gives tiny $z$, no proposal. \texttt{KELP}: ADF-$p
= 0.005$ and $|z| > 1.5$ (stated premise); fires, route to
\texttt{OU}. \texttt{SQUID\_INK}: ADF-$p = 0.41 \gg 0.05$; does not
fire.

\paragraph{Guard 4 (TSMOM).} For \texttt{SQUID\_INK}: $r^{(12\mathrm{m})}
= +1.2\%$, Olivia flag on; with
$\theta_{\text{trend}} = 0.5\%$ the guard fires, route to TSMOM
long. \texttt{RESIN}'s $+0.01\%$ is far below.

\paragraph{Guard 5 (gamma).} No options on remaining products; does
not fire.

\paragraph{Guard 6 (market making).} \texttt{RESIN}: check $\hat{s}
> 2\hat{\sigma}/\sqrt{N_{\text{fills}}}$. With $\hat{s} = 0.5$,
$\hat{\sigma} = 1.8$~bps, $N_{\text{fills}} = 100$, RHS $= 0.36$
bps $< 0.5$; depth $40 > $ min; fires, route to \texttt{MM}.

\paragraph{Guard 7 (skip).} All four products routed.

\begin{table}[htbp]
\centering
\begin{tabular}{lll}
\toprule
Product & Guard fired & Strategy \\
\midrule
\texttt{BASKET1} & Guard 2 (basket arb) & \texttt{BASKET} short \\
\texttt{KELP}    & Guard 3 (OU mean rev) & \texttt{OU} (sign from $z$) \\
\texttt{SQUID\_INK} & Guard 4 (TSMOM)    & \texttt{TSMOM} long \\
\texttt{RESIN}   & Guard 6 (market making) & \texttt{MM} \\
\bottomrule
\end{tabular}
\caption{Decision-tree output for the toy universe of
\cref{sec:wsw_toy_example}, computed guard-by-guard in
\cref{subsec:wsw_guard_eval}. Four products routed to four
different strategies by six threshold tests.}
\label{tab:wsw_toy_selection}
\end{table}

The formal output matches the intuitive assignment made at the
start of the section. The formalisation just makes it reproducible,
auditable, and tunable.

\begin{intuitionbox}
The decision tree is not a classifier but a priority-ordered
list of sufficient conditions, an if-elif chain in which each
branch's expected edge dominates the lower branches when its guard
fires. The default branch is ``do nothing'', not ``pick the most
confident other option''. Doing nothing is a strategy too, and
often the right one.
\end{intuitionbox}

% =====================================================================
\section{Meta-labeling for trade-level selection}
\label{sec:wsw_metalabel}

Not every trade the tree proposes should fire. Meta-labeling
(\cref{ch:ml_trading}) is a binary classifier that predicts whether
a proposed trade hits its profit-taker before its stop-loss. The
primary signal proposes; the meta-classifier decides.

Features: observable vector $o_t^{(p)}$, HMM regime posterior, and
trade characteristics (strategy, size, horizon). Output: probability
of profit-taker hit. Trades below threshold (typically $0.5$--$0.6$)
are dropped.

\begin{keyfactbox}[Meta-labeling in the decision framework]
Let $x_t$ be the feature vector (observables, regime posterior,
proposed-trade features) and let $y_t \in \{0, 1\}$ be the triple-barrier
label (\cref{ch:ml_trading}) of the proposed trade: $1$ if
profit-taker hit first, $0$ otherwise. Train a binary classifier
$\hat{p}(x_t) = \Prob(y_t = 1 \mid x_t)$ on historical proposals
under CPCV (\cref{ch:backtesting}). At decision time, fire the
proposed trade only if $\hat{p}(x_t) > \tau_{\text{meta}}$, where
$\tau_{\text{meta}}$ is tuned on validation data.

Typical impact: $5$--$15\%$ uplift in annualised strategy return
and $30$--$50\%$ reduction in maximum drawdown on the same
primary signal, because the meta-filter systematically drops the
trades that were about to be losers
\citeAFML{}.
\end{keyfactbox}

Meta-labeling is effective here because the primary signal is
tuned on its own optimal conditions; the tree may fire it in
slightly-off conditions where the guard passes but the edge is
thin. A regime-aware meta-classifier learns those conditions and
refuses. \Cref{ch:ml_trading} has the full worked example.

A stylised lift sketch. Primary signal (OU $z$-entry) proposes
$1000$ trades; $550$ are profitable (55\%), averaging $+10$~bps on
winners and $-8$~bps on losers. Per-trade P\&L: $0.55 \cdot 10 -
0.45 \cdot 8 = 1.9$~bps. A meta-classifier at AUC $0.60$ with
$\tau_{\text{meta}} = 0.55$ keeps $700$ trades, of which $420$ win
(60\%). Per-trade P\&L: $0.6 \cdot 10 - 0.4 \cdot 8 = 2.8$~bps,
up 47\%, with 30\% fewer trades and proportionally less turnover
cost (where the Sharpe uplift lives). Drawdown reduction is larger
still because the classifier disproportionately drops runs of
correlated losers. The shape (modest hit-rate uplift,
disproportionate Sharpe/drawdown improvement) is what \citeAFML{}
reports.

\begin{pitfallbox}[Meta-labeling is not a return predictor]
Do not train the meta-classifier to predict trade returns. Train it
to predict whether the trade hits profit-taker before stop-loss (a
binary triple-barrier label). Predicting returns is high-variance
regression that overfits; predicting barrier-hit is low-variance
classification with a stable decision boundary. \citeAFML{} is
dogmatic on this point.
\end{pitfallbox}

% =====================================================================
\section{From one strategy per product to a portfolio of strategies}
\label{sec:wsw_portfolio}

The tree so far outputs at most one strategy per product per tick.
In practice books run many strategies in parallel (market making
slow, mean reversion fast, basket arb on-dislocation). How to size
and combine them?

\textbf{Disjoint}: each strategy trades a separate slice with its
own capital and risk budget. The Frankfurt Hedgehogs did this with
seven strategies on fixed subsets. Disjoint is the simplest thing
that works and fails gracefully: one bad strategy does not
contaminate the others.

\textbf{Risk-budget combination}: allocate capital inversely to
drawdown, proportionally to Sharpe:
\begin{equation}
  w_i \;\propto\; \frac{\text{Sharpe}_i}{\text{drawdown}_i},
  \qquad \sum_i w_i = 1,
  \label{eq:wsw_risk_budget}
\end{equation}
renormalised over actively-selected strategies. Blunt, but
essentially fractional Kelly with diagonal covariance, and exactly
what the Frankfurt Hedgehogs used.

\textbf{Joint mean-variance}: estimate the covariance $\Sigma$ of
strategy returns on held-out data and solve
\begin{equation}
  \max_{\mathbf{w}\ge 0}\;
  \mathbf{w}^{\top}\boldsymbol{\mu}
  \;-\;\tfrac{\lambda}{2}\,
  \mathbf{w}^{\top}\Sigma\mathbf{w},
  \label{eq:wsw_mv_combine}
\end{equation}
which is Markowitz over strategies, not assets, with $\lambda$ the risk
aversion. $\Sigma$ is unstable on small samples; shrink or
regularise before use.

\begin{pitfallbox}[Correlated strategies are one strategy in disguise]
The expensive mistake is running two strategies that look different
but bet the same way. TSMOM and a momentum-feature meta-labeler are
near-collinear: same times, same sign, same products. Disjoint or
risk-budget combinations that treat them as independent double down
and claim false diversification. Estimate the strategy-return
correlation on held-out data and refuse pairs above $\sim 0.7$;
\citeChanQT{} uses $0.5$ as a conservative threshold.
\end{pitfallbox}

% =====================================================================
\section{Parameter adaptation: Chan's CPO}
\label{sec:wsw_cpo}

Every threshold in the decision tree
($z_{\text{entry}}$, $\theta_{\text{trend}}$,
$\tau_{\text{meta}}$, the basket-spread level, the MM depth cut,
the MM inventory skew) was written as a fixed number. Their optima
drift with conditions: a calm-regime $z$-entry is too tight for a
volatile one; a TSMOM threshold is too low in chop. Fixed
parameters across regimes is the biggest slow-bleed source
\citeChanQT{}.

The naive fix, refitting on a rolling window, has three failure
modes. \textbf{Lag}: a 60-day window cannot adapt faster than
$\sim 20$ days (how long new data needs to crowd out old), and
regimes transition faster. \textbf{Loss of significance}: shorter
windows have wider confidence intervals. \textbf{No regime
conditioning}: a rolling mean averages across whatever mix of
regimes the window contains, not the conditional-on-current-regime
optimum.

\textbf{Conditional Parameter Optimisation} (CPO)
\citep{chan2020cpo} reframes parameter optimisation as regression:
given market features $\phi_t$ (vol, HMM posterior, spread,
term-structure slope) and candidate parameter vector $\theta$,
predict held-out performance $R$. Fit $\hat{R}(\phi, \theta)$
(Chan uses gradient-boosted trees) on a historical dataset of
$\big(\phi_t, \theta, R(\phi_t, \theta)\big)$ tuples, evaluated
under CPCV. At decision time, observe $\phi_{t^*}$ and solve
\begin{equation}
  \theta^{*}(\phi_{t^*})
  \;=\;
  \arg\max_{\theta \in \Theta}\,
  \hat{R}(\phi_{t^*}, \theta).
  \label{eq:wsw_cpo}
\end{equation}
The output is the conditional optimum given current features, not
the marginal optimum averaged over history.

\begin{keyfactbox}[CPO procedure]
\begin{enumerate}
  \item \textbf{Collect training data} as tuples
  $(\phi_t, \theta, R(\phi_t, \theta))$ sampled over the
  historical period. For each historical time, sweep over a
  grid of candidate $\theta$ and evaluate held-out performance
  under CPCV.
  \item \textbf{Fit a regression} $\hat{R}(\phi, \theta)$
  (typically gradient-boosted trees) on the collected
  tuples. The regression learns the joint dependence on
  features and parameters, so it implicitly learns the
  regime-conditional optimum.
  \item \textbf{At each decision tick}, observe $\phi_{t^*}$
  and maximise $\hat{R}(\phi_{t^*}, \theta)$ over the
  candidate grid. Use the maximising $\theta^{*}$ as the
  strategy's parameters for the next tick.
\end{enumerate}
\end{keyfactbox}

CPO sits cleanly on top of the HMM because you do not need to
identify regimes explicitly: the features $\phi$ already encode
them jointly. In practice include both the HMM posterior (which
carries the sticky prior) and raw features (which carry
microstructure state).

\subsection{A toy CPO worked example}
\label{subsec:wsw_cpo_worked}

Consider a single parameter: OU entry threshold $z_{\text{entry}}$
on \texttt{KELP}. Strategy: if $|z_t^{(\mathrm{KELP})}| >
z_{\text{entry}}$, enter against the deviation, exit at $z = 0$.
Feature vector $\phi_t$ is scalar: $p_t \defeq \Prob(S_t = 2 \mid
o_{1:t})$, the HMM's high-vol posterior.

Non-CPO: sweep $z_{\text{entry}} \in \{1.0, 1.5, 2.0, 2.5, 3.0\}$,
pick the best full-history Sharpe. Suppose $z = 1.5$ wins at Sharpe
$1.70$; use $1.5$ forever.

CPO: for every historical day, record $(p_t, z_{\text{entry}},
\text{Sharpe over next 20 days})$ for each grid value. Fit
$\hat{R}(p, z_{\text{entry}})$ under CPCV. At decision time,
maximise over the grid at $p_{t^*}$. Suppose the fitted rule:
\begin{equation}
  \hat{R}(p, z_{\text{entry}})
  \;\approx\;
  \begin{cases}
    \text{peaks at } z_{\text{entry}} = 1.5 & \text{when } p \lesssim 0.3, \\
    \text{peaks at } z_{\text{entry}} = 2.0 & \text{when } 0.3 \lesssim p \lesssim 0.7, \\
    \text{peaks at } z_{\text{entry}} = 2.5 & \text{when } p \gtrsim 0.7.
  \end{cases}
  \label{eq:wsw_cpo_rule}
\end{equation}
The optimal threshold is larger in high-vol regimes because the
OU process's unconditional $\sigma$ is larger, so $z = 1.5$ is
weaker evidence of a true deviation. The non-CPO quant averages to
$1.5$; the CPO quant picks $1.5$ in calm and $2.5$ in volatile,
capturing higher OOS Sharpe in both.

Real CPO sweeps many features and parameters, but the three
ingredients are the same: $(\phi, \theta, R)$ tuples, regression
$\hat{R}$, decision $\arg\max_\theta \hat{R}(\phi_{t^*}, \theta)$.

\begin{pitfallbox}[CPO can overfit just like any supervised ML model]
CPO is supervised ML; every failure mode of \cref{ch:ml_trading}
applies. A too-flexible model class memorises noise in the
$(\phi, \theta, R)$ tuples: beautiful in-sample, random
out-of-sample. Cures: CPCV (\cref{ch:backtesting}), modest capacity
(Chan uses shallow boosted trees), and OOS-Sharpe-uplift over a
single-threshold baseline as the acceptance criterion.
\end{pitfallbox}

\begin{historybox}
CPO was popularised by Ernest Chan in the 2020s via PredictNow and
\citep{chan2020cpo}. The idea, context-aware parameter tuning,
appears in the 1980s adaptive-control literature, but Chan
packaged it for trading. Before CPO the practitioner habit was
rolling-window refit.
\end{historybox}

% =====================================================================
\section{The full pipeline}
\label{sec:wsw_pipeline}

Assemble everything. Eight steps per tick, output is a list of
sized, scheduled trades.

\begin{keyfactbox}[The full decision pipeline]
\begin{enumerate}
  \item \textbf{Observe.} Compute the observable vector
  $o_t^{(p)}$ of \cref{eq:wsw_obs_vec} for each product $p$
  in the universe, from the last $N$ ticks of order-book
  data.
  \item \textbf{Detect regime.} Run the trained HMM forward
  algorithm on the recent return history; take the last
  column as the current regime posterior.
  \item \textbf{Select strategy.} Apply the decision tree of
  \cref{sec:wsw_decision_tree} to each product. Each product
  is routed to at most one strategy or skipped.
  \item \textbf{Adapt parameters.} For each selected
  (product, strategy) pair, apply CPO
  (\cref{sec:wsw_cpo}) to pick parameters $\theta^{*}$
  conditional on the current regime and features.
  \item \textbf{Propose trades.} Run each selected strategy
  with its CPO-chosen parameters on its current observables;
  collect the proposed trades.
  \item \textbf{Meta-label filter.} For each proposed trade,
  evaluate the meta-classifier on the trade's features and
  drop trades with $\hat{p} < \tau_{\text{meta}}$.
  \item \textbf{Size.} Apply Kelly sizing
  (\cref{ch:kelly_risk}) to each surviving trade, with
  fractional Kelly ($f^{*}/4$ or $f^{*}/2$) as the
  production-safe default.
  \item \textbf{Schedule execution.} Send each sized trade
  through the Almgren--Chriss optimal execution
  trajectory (\cref{ch:almgren_chriss}). Large trades are
  broken into child orders; small trades fire marketable in
  a single shot.
\end{enumerate}
\end{keyfactbox}

Each step has a chapter behind it: step~1
(\cref{ch:lob,ch:order_flow}); step~2 (\cref{sec:wsw_hmm});
step~3 (\cref{sec:wsw_decision_tree}); step~4
(\cref{sec:wsw_cpo}); step~5
(\cref{ch:market_making,ch:mean_reversion_ou,ch:pairs_cointegration,ch:basket_etf_arb,ch:momentum,ch:arbitrage_zoo,ch:greeks_hedging});
step~6 (\cref{ch:ml_trading}); step~7 (\cref{ch:kelly_risk});
step~8 (\cref{ch:almgren_chriss}).

Validation is the domain of \cref{ch:backtesting}: each component
walk-forward or CPCV validated individually, composition validated
end-to-end to catch cross-layer leakage (e.g., using the
meta-classifier's OOS predictions as CPO input is a subtle leak).

Three implementation details. \textbf{Order is load-bearing}: CPO
(4) must follow selection (3) because the candidate grid depends
on the strategy; meta-labeling (6) must follow CPO (4) because the
classifier is trained on regime-conditional proposals; deploying
it downstream of CPO after fixed-parameter training silently
miscalibrates it (training distribution $\neq$ deployment
distribution).

\textbf{Log everything timestamped and replayable}: HMM posterior,
tree route, CPO parameters, meta-score, Kelly size, AC schedule.
Without these, debugging a live failure is impossible.

\textbf{Decouple retraining from decision cadence}: the pipeline
decides per-tick; HMM, CPO regression, and meta-classifier retrain
nightly or weekly. Retraining inside the decision loop is both a
leakage vector and a performance sink. Cron retraining writes
versioned artefacts; the decision loop loads the latest at a
barrier.

\begin{intuitionbox}
The pipeline is a funnel. Regime detection narrows ``any strategy''
to ``regime-appropriate''; the tree narrows to one per product; CPO
picks parameters; strategies propose; meta-labeling drops the
worst; Kelly sizes; AC schedules. Every step throws something
away; nothing is added. By the time a proposal becomes an order,
every filter has had a veto.
\end{intuitionbox}

\begin{pitfallbox}[The pipeline's weakest step is usually the HMM]
Step 2 is usually the weak link. A Gaussian HMM on a short window,
wrong number of states, or non-Gaussian tails produces regime
labels wrong for half of history, and every downstream layer
conditions on them. Monitor regime detection in isolation: plot
the forward posterior against crashes/VIX spikes, watch
transition-matrix stickiness for drift, retrain slowly enough that
the sticky prior does not whipsaw. Never bury the HMM where you
cannot see its output.
\end{pitfallbox}

\begin{gsbox}
On a Goldman central-risk desk, a slimmer pipeline runs
continuously across the single-stock book. Strat owns the pipeline
code (regime model, tree, CPO regression, meta-classifier). The
trader owns the override, the red button that disables the
regime layer for N ticks or forces strategy $X$ on product $Y$. The
quant researcher owns the regime model's retraining cadence and
investigates anomalies. When it works the pipeline runs for weeks
untouched; when regime detection breaks, all three scramble to the
same stand-up. This Strat/trader/researcher triad is textbook at
every tier-1 bank \citeChanQT{}.
\end{gsbox}

\begin{historybox}
Regime-conditional strategy selection appears in the practitioner
literature of the 2000s; Chan's \emph{Quantitative Trading} (2009,
\citeChanQT{}) was the first retail-facing book to present it with
code. Academic HMM-for-regime papers (Hamilton's 1989
Markov-switching model) predate it by two decades. Meta-labeling
was popularised by L\'opez de Prado (2018, \citeAFML{}), ch.\ 3.
CPO is Chan's more recent contribution \citep{chan2020cpo}. The
pipeline assembled here is a mid-2020s synthesis: components
all at least a decade old, but rarely written down combined.
\end{historybox}

% =====================================================================
\section{Alternatives and honest limitations}
\label{sec:wsw_alternatives}

Three alternatives to the pipeline are worth knowing.

\paragraph{Always-on ensemble.} Drop the regime layer; run every
strategy all the time with a vol-targeted ensemble. Advantage: no
missed trades from HMM confusion. Disadvantage: every tick eats
every strategy's worst-regime loss. Backtests usually favour the
pipeline when the HMM is well-calibrated (sign flips are sharp),
and the always-on ensemble when the HMM is badly calibrated. Start
always-on, measure regime sensitivity, gate with the HMM only once
it beats straight averaging.

\paragraph{Black-box end-to-end model.} Train one large model (NN,
RL, transformer) to emit trade proposals directly. Appeal: no
hand-coded tree, no HMM, no CPO. Reality: every quant RL group has
tried it and produced Sharpe-$<$-1 models (\cref{ch:ml_trading};
financial series are not stationary and do not contain enough
information to train end-to-end). The layered pipeline works
\emph{because} each layer is small and imposes structure (HMM:
Markovian regimes; tree: priority; CPO: conditional regression;
meta: triple-barrier), and each can be validated independently. An
end-to-end model is a monolithic failure mode.

\paragraph{Manual discretion.} A human reads the screen and picks a
strategy, implicitly running the tree. Advantage: qualitative
state (news, CB announcement, block cross) enters the decision.
Disadvantage: slow, unreproducible, cannot scan $10{,}000$
products. The compromise every desk converges on is pipeline +
trader override (as in \cref{sec:wsw_pipeline}): pipeline does the
reproducible part; human handles the rest.

Limitations worth flagging. The pipeline assumes regimes exist
and are detectable (not universally true); assumes the strategy
pool covers plausible states (novel regimes such as pandemic, war,
or rate shocks may fit nothing); assumes clean input data (bad
ticks and outages need their own pre-processing layer, not
discussed here). Every production desk has a dedicated monitoring
layer.

% =====================================================================
\section{A short reference implementation}
\label{sec:wsw_code}

The pipeline is a pure function. We show only the glue, not the
underlying HMM, CPO, meta-classifier, or strategy implementations;
each object exposes the interface of its own chapter.

\begin{lstlisting}[language=Python, caption={Pipeline composition.
Each dependency object implements the interface of its chapter.
See \cref{sec:wsw_pipeline} for the eight-step breakdown.}]
def decide_trades(universe, market_state, hmm, decision_tree,
                  cpo_model, meta_clf, kelly, execution):
    """Return a list of sized, scheduled child orders."""
    # Step 2: regime posterior on most recent return history
    regime = hmm.forward_posterior(market_state.recent_returns())
    trades = []
    for product in universe:
        # Step 1: observables
        obs = market_state.observe(product)
        # Step 3: strategy selection
        strategy = decision_tree.select(obs, regime)
        if strategy is None:
            continue
        # Step 4: CPO parameter adaptation
        params = cpo_model.predict_params(strategy.name, obs, regime)
        # Step 5: primary proposal
        proposed = strategy.propose(obs, params)
        if proposed is None:
            continue
        # Step 6: meta-label filter
        features = meta_clf.build_features(obs, regime, proposed)
        if meta_clf.predict_proba(features) < 0.5:
            continue
        # Step 7: Kelly sizing
        sized = kelly.size(proposed, regime=regime)
        # Step 8: Almgren--Chriss child-order schedule
        children = execution.schedule(sized)
        trades.extend(children)
    return trades
\end{lstlisting}

Thirty lines of glue over eight chapters of machinery. Reading
\texttt{decide\_trades} with Part~V loaded is a walk through the
funnel, not new material.

% =====================================================================
\section{Key facts to memorise}
\label{sec:wsw_keyfacts}

\begin{itemize}
  \item The chapter picks which strategy to run: observable in,
  sized trade list out.
  \item Regime detection with a 2-state Gaussian HMM: forward
  algorithm for the current posterior, Viterbi for historical
  paths, Baum--Welch for parameters. $N = 2$ usually; more
  overfits.
  \item Transition-matrix stickiness sets flip-sensitivity:
  diagonal $0.98$ gives $50$-day mean dwell time. A single
  $-2.5\%$ return flipped the worked example from $0.91 \to
  0.21$ on state 1 in one step.
  \item Regime-to-strategy map: mean reversion wins in
  range-bound regimes, momentum in trending, gamma in high-vol,
  arbitrage/basket are regime-neutral, market making wants low
  adverse selection (calm or range-bound).
  \item Synthetic: TSMOM Sharpe $+1.65$ trending / $-1.97$
  mean-reverting; OU exactly swapped ($-1.64$ / $+2.22$). Wrong
  regime is a sign flip, not a small penalty.
  \item Decision tree priority: arbitrage, basket, OU/pairs,
  TSMOM, gamma, MM, SKIP. Each guard is a threshold on one
  component of the observable vector.
  \item Meta-labeling: binary classifier on proposed trades
  trained on triple-barrier labels. Typical impact $5$--$15\%$
  Sharpe uplift, $30$--$50\%$ drawdown reduction. Predict
  barrier-hit, never returns.
  \item CPO: parameter adaptation as regression on $(\phi,
  \theta, R)$ tuples. Sits on top of HMM posterior; replaces
  rolling-window refit.
  \item Pipeline: observe, detect, select, adapt, propose,
  filter, size, schedule. Every step throws trades away;
  nothing is added.
  \item At GS (and every tier-1 bank): Strat owns code, trader
  owns override, quant researcher owns regime model.
  \item Biggest failure mode: unmonitored HMM. Wrong regime
  label corrupts everything downstream. Monitor the posterior
  against events; retrain slower than the sticky prior.
\end{itemize}
